% =========================================================================
% Chapter 9 - The training pipeline
% =========================================================================
\chapter{The training pipeline}
\label{ch:pipeline}

\section{Requirements}
\label{sec:requirements}

\begin{pnote}
Poraqu\^e requires \textbf{Python 3.11 or newer}
(\opt{requires-python = ">=3.11"}). Every module has been verified against the
3.11 grammar, so nothing in the source uses syntax or standard-library calls
newer than that. There is deliberately no upper bound: newer interpreters are
supported.
\end{pnote}

\begin{lstlisting}
git clone https://github.com/seixas-research/poraque.git
cd poraque
pip install -e .
\end{lstlisting}

This pulls in NumPy, SciPy, ASE, Matplotlib, PyYAML and PyTorch. A working
\opt{pdflatex} or \opt{latexmk} is needed only for the automatic PDF reports
of Section~\ref{sec:pdfreports}; without one the pipeline still runs and simply
writes the \LaTeX{} source instead.

\section{Configuration}

A run is defined by a YAML file with four sections --- \opt{data},
\opt{model}, \opt{training}, \opt{output} --- mirroring the four things a run
needs. Command-line flags override individual entries, so one committed
configuration can be swept from the shell without being edited or copied:

\begin{lstlisting}
poraque-train --write-config configs/train_config.yaml
poraque-train --config configs/train_config.yaml --epochs 500
\end{lstlisting}

Unknown keys are \emph{rejected} rather than ignored: a typo such as
\opt{learn\_rate} would otherwise be silently dropped and the run would quietly
use the default, producing results that do not match the file that appears to
describe them. The \emph{resolved} configuration is written beside the results,
recording what actually ran, overrides included.

\section{The external potential is computed, not imported}
\label{sec:computedextcar}

Poraqu\^e \textbf{computes} the external potential natively, from
\file{POSCAR}, \file{INCAR} and \file{POTCAR}, using the tabulated
pseudopotential of Chapter~\ref{ch:vasp}. There is no option to import one:
the training input must be exactly what the pipeline produces at inference
time, when no such file exists.

This costs essentially nothing in fidelity. Against reference potentials on
five gold structures:

\begin{center}
\begin{tabular}{@{}lrrr@{}}
\toprule
Structure & grid & relative $L^2$ & Pearson $r$ \\
\midrule
struct\_000 & $128^3$ & $2.13\times10^{-5}$ & $1.000000$ \\
struct\_001 & $120{\times}128{\times}128$ & $5.95\times10^{-5}$ & $1.000000$ \\
struct\_002 & $128{\times}120{\times}128$ & $5.75\times10^{-5}$ & $1.000000$ \\
struct\_003 & $128^3$ & $5.66\times10^{-5}$ & $1.000000$ \\
struct\_004 & $128^3$ & $5.36\times10^{-5}$ & $1.000000$ \\
\bottomrule
\end{tabular}
\end{center}

\begin{ptip}
The shared grid is taken from \file{CHGCAR}, not \file{EXTCAR}: the density is
present in any standard run, the external potential is not. Sourcing the grid
from a file that may not exist would have reintroduced the dependency through
the back door.
\end{ptip}

\section{Gaussian smoothing of the potential}
\label{sec:smoothing}

The tabulated pseudopotential is sharply peaked at the ionic cores. Since a
network with a fixed mode truncation may not use the highest spatial
frequencies, an optional Gaussian blur is available:

\begin{lstlisting}
poraque-train --config configs/train_config.yaml \
       --gaussian-blur 0.15 --blur-method spectral
\end{lstlisting}

Smoothing is applied \emph{on top of} the tabulated potential; it never
replaces it. It is off by default (\opt{gaussian\_blur: null}).

\subsection{Two methods, and why the default matters}

\begin{description}
  \item[\opt{spectral} (default)] multiplies by $e^{-G^2\sigma^2/2}$ in
    reciprocal space. Exact for a band-limited field, and it uses the true
    reciprocal metric, so the blur stays isotropic in Cartesian space on
    \emph{any} cell.
  \item[\opt{ndimage}] uses \opt{scipy.ndimage.gaussian\_filter} with
    \opt{mode="wrap"} --- the wrap is essential, since the cell is a torus.
\end{description}

\begin{pwarn}
The two agree to a relative $1.4\times10^{-5}$ on an \emph{orthogonal} cell and
differ by $2$--$6\,\%$ on the fcc cell of the reference data. \opt{ndimage}
blurs along \emph{lattice} axes, which is anisotropic in Cartesian space
whenever the lattice vectors are not orthogonal. For a face-centred cubic
lattice --- angles of $60^\circ$ --- it is measurably wrong, not merely
different.
\end{pwarn}

\subsection{Measured effect}

The blur is far from uniform. Binned by distance to the nearest ion at
$\sigma = 0.15$~\AA:

\begin{center}
\begin{tabular}{@{}lrrrr@{}}
\toprule
Shell [\AA] & voxels & mean $|\Delta V|$ [eV] & $\overline{V}$ [eV] & $\overline{\rho}$ [e/\AA$^3$] \\
\midrule
$0$--$0.5$ & $959$ & $9.64$ & $-113.8$ & $3.51$ \\
$0.5$--$1.0$ & $6748$ & $4.06$ & $-45.8$ & $1.46$ \\
$1.0$--$1.5$ & $17854$ & $1.58$ & $+13.0$ & $0.32$ \\
$1.5$--$2.0$ & $7194$ & $1.18$ & $+25.8$ & $0.16$ \\
\bottomrule
\end{tabular}
\end{center}

The core region ($< 1$~\AA) is $23.5\,\%$ of the voxels but absorbs
$50\,\%$ of the total change, and the deepest well shrinks by $7.8\,\%$
($-137.3 \to -126.6$~eV). Those same voxels carry the charge-density peaks
($\overline\rho = 1.72$ against a cell mean of $0.61$). Smoothing therefore
preferentially removes exactly the information that generates the density
maxima.

\begin{pnote}
On a matched held-out comparison --- same four training structures, same seed,
architecture and epochs, sole occupancy of the accelerator --- blurring changed
the held-out relative $L^2$ from $0.0274$ to $0.0272$, i.e.\ by $0.7\,\%$, and
training time by $0.4\,\%$. \textbf{Neither is a meaningful difference.} At
$32^3$ the field is already band-limited at $|G|_{\max}\approx26\ \text{\AA}^{-1}$
and the blur removes content the mode truncation was discarding anyway. The
option is worth revisiting only at higher resolution or larger \opt{modes},
where the model could actually use those frequencies.
\end{pnote}

\section{Training}
\label{sec:universal}

\textbf{One model per task is trained across all the training structures},
never one per material.

\begin{lstlisting}
poraque-train --config configs/train_config.yaml
\end{lstlisting}

There is one protocol and one variation on it: \opt{valid\_fraction} holds back
a fifth of the structures by default, and \opt{enable\_kfold}
(Section~\ref{sec:kfold}) swaps the whole thing for cross-validation. Nothing
else selects a training mode.

Batches are drawn \emph{across} materials: the sampler groups structures by
grid shape and shuffles both within and across those groups, so a gradient step
generally mixes several. On the reference data the buckets are
$(32,32,32)\times3$, $(30,32,32)\times1$ and $(32,30,32)\times1$ --- three
distinct shapes flowing through one set of weights.

The run writes a single unified checkpoint,
\file{models/poraque\_models.pfno}, holding both operators under the keys
\opt{ext2chg} and \opt{chg2tau}.

\begin{pwarn}
With \opt{valid\_fraction: 0} every structure is in training, so the reported
metrics are \textbf{training fit}. They show the model can represent the data;
they are not a generalisation estimate. Raise \opt{valid\_fraction}, or use the
cross-validation of Section~\ref{sec:kfold}, for that.

Measured on seventeen gold structures: the all-structure training fit gives
relative $L^2$ of $0.0209$ (\opt{ext2chg}) and $0.0140$ (\opt{chg2tau}),
against held-out values of $0.0379$ and $0.0511$ --- roughly two to four times
larger. That gap is the memorisation margin, and quoting the training-fit
number as a result would be a serious overstatement.
\end{pwarn}

\section{K-fold cross-validation}
\label{sec:kfold}

\begin{lstlisting}
poraque-train --config configs/train_config.yaml \
       --kfold --k-folds 5
\end{lstlisting}

or in the configuration file:

\begin{lstlisting}
training:
  enable_kfold: true
  k_folds: 5
\end{lstlisting}

Each fold trains a fresh model on $K-1$ groups and scores it on the held-out
group, in physical units. The result is a generalisation estimate \emph{with a
spread}, not a model to deploy.

\subsection{The split is at the structure level}

This is the part that decides whether the number means anything.

A fold is a set of \emph{material names}; every voxel of a material goes to the
same side of the split. Splitting at the voxel level would place the same
material in both training and validation, and since neighbouring voxels of one
crystal are enormously correlated, the model would effectively be scored on
data it had already seen. The resulting figure would look excellent and would
say nothing about transfer to a new material.

\begin{ptip}
\opt{k\_folds} is capped at the number of structures; setting it equal to that
count \emph{is} leave-one-out. The partition is verified exact and disjoint ---
every structure appears in exactly one validation fold.
\end{ptip}

\section{Automatic PDF reporting}
\label{sec:pdfreports}

Every run assembles its metrics and figures into a typeset PDF under
\file{reports/}. In cross-validation mode a single \emph{consolidated} report
covers all folds, with the per-fold table and the mean $\pm$ standard deviation
of each metric.

The mechanism matters for hygiene: the \LaTeX{} source is written into a
\opt{tempfile.mkdtemp()} scratch directory, compiled there, and only the PDF is
moved out; the whole tree is then removed with \opt{shutil.rmtree}.

\begin{ptip}
Deleting the directory wholesale is the only reliable way to guarantee that no
\file{.tex}, \file{.aux}, \file{.log}, \file{.toc}, \file{.fls} or
\file{.fdb\_latexmk} survives --- a hand-maintained list of extensions would
drift as packages begin emitting new ones. Compilation falls back from
\opt{latexmk} to two \opt{pdflatex} passes, and if no toolchain is present the
source is written beside the target and the caller is told, rather than failing
silently.
\end{ptip}

Each report carries a \emph{What these numbers do not show} section, generated
from the run itself: whether anything was held out, how many structures there
were, and whether the figures support a generalisation claim.

\section{Outputs}

\begin{center}
\begin{tabular}{@{}ll@{}}
\toprule
Path & Contents \\
\midrule
\file{models/poraque\_models.pfno} & both trained operators, in one file \\
\file{logs/*.log} & human-readable training log \\
\file{logs/*.json} & machine-readable metrics and full history \\
\file{logs/*\_config.yaml} & the resolved configuration, for reproduction \\
\file{results/plots/} & loss curves, cross-sections, parity plots \\
\file{reports/*.pdf} & typeset report per model or per cross-validation \\
\bottomrule
\end{tabular}
\end{center}
